\documentclass[11pt,a4paper]{appletmanual}
\usepackage[spanish,es-noquoting,es-nodecimaldot]{babel}
\manuallang{es}

\appletname{Beneficios, Costos y Oferta}
\appletsub{El problema de la empresa competitiva en el corto plazo, y la curva
de oferta que sale de él. Dos paneles dibujados a la misma escala, porque son
la misma curva.}
\appletdir{firm-supply}

\begin{document}
\manualtitlepage{Manual del applet · Microeconomía I · Universidad de los Andes\\
Traducido del applet de GeoGebra «Beneficios y Costos».\\
Código MIT. Se abre en el navegador; no hay nada que instalar.}

\manualtoc
\thispagestyle{empty}
\clearpage
\setcounter{page}{1}

% ==========================================================================
\section{Qué es esto}

Una empresa competitiva en el corto plazo tiene un solo grado de libertad:
cuánto producir. El capital está fijo, el precio se lo dan hecho, y todo lo que
puede hacer es elegir la cantidad $q$ que maximiza el beneficio
\[
  \pi(q) \;=\; p\,q \;-\; c(q).
\]

Ese es el problema. La curva de oferta de la empresa es la \emph{respuesta} a
ese problema, resuelta una vez por cada precio posible. Son dos objetos
distintos, y en casi todos los manuales aparecen en dos figuras distintas y en
dos páginas distintas, con una frase que dice que el costo marginal «es» la
curva de oferta y que el estudiante tiene que creerse.

Este applet los pone uno al lado del otro y a la misma escala. El panel derecho
es el problema; el izquierdo es su respuesta. El costo marginal se dibuja en
los dos, en las mismas coordenadas, y en el panel de oferta va superpuesto en
rojo discontinuo sobre la curva de oferta. Por encima del precio de cierre las
rayas caen exactamente encima. No hay que creerse nada.

\figher[1.0]{panels}{Los dos paneles con los ajustes de partida. Izquierda: la
curva de oferta de la empresa en el plano $(q, p)$. Derecha: el problema, con
el costo marginal, el medio y el variable medio. Los dos ejes son los mismos en
los dos paneles, así que el costo marginal ocupa el mismo lugar en ambos.}

\begin{amkey}
El costo marginal y la oferta inversa no son curvas parecidas: por encima del
mínimo del costo variable medio son la misma curva, punto por punto. Todo el
applet está construido para que eso se vea en vez de enunciarse.
\end{amkey}

\subsection{Para quién}

Para un primer curso de microeconomía, en la sesión en que se pasa de la
maximización de beneficios a la curva de oferta. Supone que el estudiante ya ha
visto costos totales, medios y marginales, y que sabe qué es un óptimo. No
supone cálculo más allá de la idea de pendiente.

%% Secciones comunes a todos los manuales, edición española.
%% Se incluyen con \input{common-es}. Lo único que cambia entre applets es
%% \appletfolder, que es también el nombre de la carpeta y de la ruta en el sitio.

\section{Cómo abrirlo}

No hay que instalar nada. La aplicación entera es un solo archivo
\code{index.html}: no descarga librerías, no llama a ningún servidor y no
guarda nada fuera del navegador.

\subsection{Las tres maneras}

\begin{description}
\item[Doble clic sobre el archivo.] Es lo más rápido. Busque
  \code{\appletfolder/index.html} y ábralo. El navegador lo muestra desde el disco,
  con la barra de direcciones empezando por \code{file://}. Funciona sin
  conexión desde el primer momento.
\item[Desde la web.] Si alguien ya lo publicó, la dirección termina en
  \code{/\appletfolder/}. Una vez cargada la página, puede desconectarse: no
  necesita la red para nada más.
\item[Servido desde una carpeta.] Para una clase, poner la carpeta en
  cualquier servidor estático basta. Con Python instalado:
  \begin{quote}\code{python3 -m http.server 8000}\end{quote}
  y luego \code{http://localhost:8000/\appletfolder/} en el navegador.
\end{description}

\begin{amnote}
Guardar la página con \key{Ctrl}\,+\,\key{S} (o \key{Cmd}\,+\,\key{S} en
Mac) deja una copia propia que sigue funcionando sin conexión y que se puede
llevar en una memoria USB o repartir por correo.
\end{amnote}

\subsection{Idioma y tema}

Arriba a la derecha hay dos pares de botones.

\begin{ctrltable}{Control & Qué hace}
ES / EN & Cambia el idioma de toda la interfaz al vuelo. Nada se recalcula:
          puede cambiarlo en mitad de una explicación sin perder el estado. \\
Oscuro / Claro & Cambia el tema. El oscuro es el diseño original; el claro
          existe para proyectores y para imprimir. \\
\end{ctrltable}

Las dos elecciones se recuerdan en el navegador, así que la próxima vez se abre
como la dejó.

\subsection{Qué navegador}

Cualquiera que sea razonablemente reciente: Chrome, Edge, Firefox o Safari, en
ordenador, tableta o teléfono. En pantallas estrechas los paneles se apilan uno
sobre otro en lugar de ponerse al lado. No hace falta cuenta, ni permisos, ni
extensiones.


% ==========================================================================
\section{La pantalla}

\fig[1.0]{interface}{La pantalla completa. A la izquierda la columna de
controles; arriba el veredicto; abajo los dos paneles con sus lecturas.}

De izquierda a derecha y de arriba abajo:

\begin{description}
\item[La columna de controles.] Elige la función de costos, mueve sus
  parámetros, fija el precio, decide qué muestra el panel de la empresa y qué
  capas se dibujan. Los deslizadores tienen a su derecha una casilla numérica:
  se puede arrastrar o escribir el número exacto.
\item[El veredicto.] La banda de color bajo el título. Nombra en qué caso está
  la empresa —beneficio, pérdida, cierre, nivelación, capacidad, sin óptimo— y
  explica por qué. Cambia solo al mover cualquier cosa.
\item[El panel de oferta] (izquierda). La curva de oferta inversa en el plano
  $(q, p)$, con el precio de cierre, el de nivelación y el precio actual.
\item[El panel de la empresa] (derecha). O bien las curvas marginales, o bien
  las totales, según el conmutador \ui{Panel de la empresa}.
\item[Las lecturas.] La fila de números al pie de cada panel. Son los mismos
  valores que están dibujados, escritos con tres decimales para cuando el ojo
  no basta.
\end{description}

\subsection{Arrastrar para cambiar el precio}

En cualquier panel cuyo eje vertical \emph{sea} el precio —el de oferta
siempre, el de la empresa cuando muestra marginales— se puede arrastrar con el
ratón o el dedo y el precio sigue al puntero. Es la forma natural de leer una
curva de oferta: se elige un precio y se mira qué cantidad sale.

En la vista de totales el eje vertical es dinero, no dinero por unidad, y el
panel no responde al puntero.

\subsection{Teclado}

\begin{ctrltable}{Tecla & Qué hace}
\key{1} \key{2} \key{3} \key{4} & Elige la función de costos. \\
\key{m} & Panel de la empresa en marginales. \\
\key{t} & Panel de la empresa en totales. \\
\end{ctrltable}

\subsection{El tema oscuro}

El diseño original es oscuro; el claro está pensado para proyectores y para
imprimir. Las dos versiones se diseñaron por separado, no es una inversión de
colores.

\fig[0.95]{dark}{Los mismos dos paneles en el tema oscuro.}

% ==========================================================================
\section{El modelo}

\subsection{Costos}

Todo parte de una función de costo total $c(q)$: lo que cuesta producir $q$
unidades cuando el capital ya está pagado. De ella salen las demás.

\begin{align*}
  \text{Costo fijo}          &&  F   &= c(0) \\
  \text{Costo variable}      &&  CV(q) &= c(q) - F \\
  \text{Costo marginal}      &&  CMg(q) &= c'(q) \\
  \text{Costo medio}         &&  CMe(q) &= \frac{c(q)}{q} \\
  \text{Costo variable medio}&&  CVMe(q) &= \frac{c(q) - F}{q}
\end{align*}

El costo fijo se lee como el límite de $c(q)$ cuando $q$ tiende a cero por la
derecha. En las familias que vienen de una función de producción ese límite es
$r\bar K$, la renta del capital fijo: se paga se produzca lo que se produzca.

\subsection{El problema}

\[
  \max_{q \,\ge\, 0} \;\; \pi(q) \;=\; p\,q - c(q).
\]

Si el máximo cae en el interior, la condición de primer orden da
$CMg(q^\ast) = p$, y la de segundo orden exige que el costo marginal esté
\emph{creciendo} ahí. Las dos cosas importan: en una curva de costos cúbica la
ecuación $CMg(q) = p$ tiene dos raíces y una de ellas es un \emph{mínimo} del
beneficio.

Pero $q^\ast = 0$ también es un candidato, y no es un caso límite del anterior:
es una esquina, con beneficio $\pi(0) = -F$. La empresa produce solo si le va
mejor produciendo que cerrando.

\subsection{El precio de cierre}

Comparando las dos opciones:
\[
  p\,q - c(q) \;>\; -F
  \quad\Longleftrightarrow\quad
  p\,q \;>\; c(q) - F = CV(q)
  \quad\Longleftrightarrow\quad
  p \;>\; CVMe(q).
\]

El costo fijo desaparece de la comparación. Ya está pagado, y una decisión que
lo tenga en cuenta es una decisión mal tomada. De ahí sale el
\textbf{precio de cierre}:
\[
  \underline{p} \;=\; \min_{q > 0} CVMe(q),
\]
y la regla que todo el mundo recita: la empresa produce mientras el precio
cubra el costo variable medio, aunque no cubra el medio total.

\subsection{El precio de nivelación}

\[
  p_{ne} \;=\; \min_{q>0} CMe(q).
\]
Es el precio al que el beneficio es exactamente cero. Por encima hay beneficio,
por debajo hay pérdidas —que pueden convenir igualmente, mientras el precio
siga por encima de $\underline{p}$.

\subsection{La oferta}

\[
  q^\ast(p) \;=\;
  \begin{cases}
    0 & \text{si } p < \underline{p},\\[2pt]
    \text{la raíz de } CMg(q) = p \text{ en la rama creciente} & \text{si } p > \underline{p}.
  \end{cases}
\]

En $p = \underline{p}$ la empresa es indiferente, y la oferta \textbf{salta}:
pasa de cero a la cantidad $\underline{q}$ que minimiza el costo variable
medio, sin pasar por los valores intermedios. Ese salto no es un detalle
técnico. Es la razón por la que la curva de oferta de una empresa con costos
fijos no llega hasta el origen.

\subsection{El excedente del productor}

\[
  EP \;=\; \int_0^{q^\ast}\!\bigl(p - CMg(q)\bigr)\,dq
       \;=\; p\,q^\ast - CV(q^\ast)
       \;=\; \pi + F.
\]

Es el mismo número escrito de tres maneras, y el applet lo dibuja de dos: en el
panel derecho como el área entre $p$ y el costo marginal, y en el izquierdo
como el área entre el eje de precios y la curva de oferta, desde
$\underline{p}$ hasta $p$. Las dos regiones tienen formas distintas —con un
salto no pueden coincidir— y exactamente la misma área.

\begin{amnote}
Excedente del productor y beneficio son lo mismo solo cuando no hay costo fijo.
Con $F > 0$ el excedente es siempre mayor, exactamente en $F$.
\end{amnote}

% ==========================================================================
\section{Las cuatro funciones de costos}

Se eligen con los cuatro botones de arriba de la columna, o con \key{1} a
\key{4}. Cada una trae sus propios deslizadores.

\subsection{1 · Cuadráticos / potencia}

\[
  c(q) = F + a\,q^{k}, \qquad k > 1.
\]
La forma de manual. El costo marginal, $CMg = a k q^{k-1}$, crece desde cero,
de modo que el costo variable medio no tiene mínimo interior: su ínfimo es cero
y se alcanza solo en el límite. Consecuencia: \textbf{esta empresa nunca
cierra}. A cualquier precio positivo le conviene producir algo.

\begin{ctrltable}{Deslizador & Qué es}
$F$ & Costo fijo. Sube o baja el costo medio sin tocar el marginal. \\
$a$ & Escala del costo variable. \\
$k$ & Curvatura. Con $k$ cerca de 1 el costo marginal es casi constante y el
      problema deja de tener solución; con $k$ grande sube muy deprisa y la
      oferta se hace vertical. \\
\end{ctrltable}

\subsection{2 · Personalizada}

Se escribe $c(q)$ en la casilla de texto. El costo marginal se obtiene
derivando la expresión de forma \emph{simbólica}, no numérica, así que es
exacto. La que viene por defecto,
\[
  c(q) = 0{,}1\,q^{3} - 0{,}5\,q^{2} + 2q,
\]
es la cúbica del manual: costo fijo nulo, costo variable medio en U con mínimo
$1{,}375$ en $q = 2{,}5$, y por tanto precio de cierre estrictamente positivo y
un salto visible en la oferta.

\fig[1.0]{cubic}{La cúbica a $p = 1{,}6$. En el panel de oferta se ven las tres
piezas de la curva: pegada al eje de precios por debajo de $1{,}375$, el salto
discontinuo, y la rama creciente que coincide con el costo marginal.}

\subsection{3 · Producción Cobb--Douglas}

En el corto plazo el capital está fijo en $\bar K$ y el único factor variable es
el trabajo. Partiendo de
\[
  f(K, l) = b\,\bigl(K^{\beta} l^{1-\beta}\bigr)^{s},
\]
se invierte para obtener el trabajo necesario para producir $q$,
\[
  l(q) = \Bigl(\frac{q}{b\,\bar K^{\beta s}}\Bigr)^{\!e}, \qquad
  e = \frac{1}{(1-\beta)s},
\]
y se valora al salario $w$:
\[
  c(q) \;=\; r\bar K \;+\; w\,b^{-e}\,\bar K^{-\beta/(1-\beta)}\,q^{\,e}.
\]

Todo depende de $e$. Con $e > 1$ hay rendimientos decrecientes del factor
variable, el costo marginal crece y el problema tiene solución. Con
$e \le 1$ el costo marginal no crece, cada unidad extra añade beneficio, y no
hay cantidad óptima. El applet lo dice en vez de dibujar algo sin sentido.

\fig[0.8]{no-optimum}{El veredicto cuando $e \le 1$: no hay óptimo, y se explica
qué mover para que lo haya.}

\begin{ctrltable}{Deslizador & Qué es}
$b$ & Productividad total. \\
$\beta$ & Elasticidad del capital. Entra en $e$. \\
$s$ & Rendimientos a escala de la Cobb--Douglas. Entra en $e$. \\
$w$ & Salario. Escala todo el costo variable. \\
$r$ & Renta del capital. Es el costo fijo, $F = r\bar K$. \\
$\bar K$ & El capital fijo. Es lo que hace que el corto plazo sea corto. \\
\end{ctrltable}

\subsection{4 · Producción log-CD logística}

La misma idea, con una función de producción que \emph{satura}:
\[
  g(K, l) = \frac{b\,K^{\gamma}}{1 + e^{-s(\ln f - m)}},
\]
cuyo techo es $\bar Q = b\,\bar K^{\gamma}$ por mucho trabajo que se contrate.
Invirtiéndola,
\[
  c(q) \;=\; r\bar K \;+\; w\,e^{m/(1-\beta)}\,\bar K^{-\beta/(1-\beta)}
             \left(\frac{q}{\bar Q - q}\right)^{\!e},
  \qquad 0 < q < \bar Q .
\]

El costo marginal se dispara al acercarse a $\bar Q$, y la curva de oferta tiene
una asíntota vertical ahí. Es la representación honesta de lo que significa
tener el capital fijo: por mucho que suba el precio, hay una cantidad que esta
planta no puede pasar.

\fig[1.0]{capacity}{La familia logística a $p = 3{,}2$. La capacidad
$\bar Q = b\bar K^{\gamma}$ está marcada con una vertical gris, y el costo
marginal sube hacia ella sin llegar nunca.}

% ==========================================================================
\section{Los controles}

\subsection{Precio}

Un deslizador de 0 a 10, y la casilla numérica para escribir un valor exacto.
Bajo él aparece una frase que dice dónde está el precio respecto del de cierre
y del de nivelación. También se puede mover arrastrando dentro de cualquier
panel cuyo eje vertical sea el precio.

\subsection{Panel de la empresa}

\begin{ctrltable}{Vista & Qué dibuja}
Marginales & $CMg$, $CMe$ y $CVMe$, en dinero por unidad. Es la misma unidad
  que el precio, así que este panel comparte \emph{los dos} ejes con el de
  oferta y las dos figuras son literalmente comparables píxel a píxel. \\
Totales & $c(q)$, el ingreso $pq$ y el beneficio $\pi(q)$, en dinero. El eje
  vertical ya no es el del panel de oferta; el de cantidad sí lo sigue siendo.
  La pastilla de la cabecera del panel dice cuál de los dos casos es. \\
\end{ctrltable}

\subsection{Las capas}

La lista cambia según la vista, porque no todas las capas tienen sentido en las
dos.

\subsubsection{En marginales}

\begin{ctrltable}{Capa & Qué dibuja}
Costo marginal & $CMg(q)$, en rojo. También aparece superpuesto en el panel de
  oferta. \\
Costo medio & $CMe(q)$, en oro. Su mínimo es el precio de nivelación. \\
Costo variable medio & $CVMe(q)$, en naranja. Su mínimo es el precio de cierre. \\
Línea de precio & La horizontal en $p$. \\
Área de costos variables & La región bajo $CMg$ hasta $q^\ast$, que es
  $CV(q^\ast)$. \\
Área de excedente & La región entre $p$ y $CMg$ hasta $q^\ast$. \\
Óptimo & El punto $(q^\ast, p)$, su vertical, y el rectángulo del beneficio
  entre $p$ y $CMe(q^\ast)$, verde si hay beneficio y rojo si hay pérdidas. \\
Punto de equilibrio & Marca los mínimos de $CMe$ y de $CVMe$. \\
Cuadrícula & La retícula de fondo. \\
\end{ctrltable}

\fig[1.0]{areas}{Las dos áreas encendidas a la vez, con la cúbica a
$p = 1{,}6$. A la izquierda el excedente medido como integral de $q$ respecto
del precio; a la derecha, el mismo número medido como integral de $p - CMg$
respecto de la cantidad.}

\subsubsection{En totales}

\begin{ctrltable}{Capa & Qué dibuja}
Costo total & $c(q)$, con una horizontal punteada en el costo fijo. \\
Ingreso total & La recta $pq$. \\
Beneficio & $\pi(q) = pq - c(q)$. Su máximo está en $q^\ast$. \\
Costo variable & $c(q) - F$, la misma curva bajada. \\
Tangente de pendiente $p$ & La recta tangente a $c(q)$ en $q^\ast$. Es paralela
  a la recta de ingresos: ahí está la condición de primer orden, dibujada. \\
Óptimo & El segmento vertical entre $c(q^\ast)$ y $pq^\ast$, cuya longitud
  \emph{es} el beneficio. \\
Punto de equilibrio & Las cantidades donde $\pi(q) = 0$. \\
\end{ctrltable}

\fig[1.0]{totals}{La vista de totales con la cúbica a $p = 1{,}6$. La tangente
punteada toca el costo justo donde su pendiente iguala el precio, y el segmento
vertical en $q^\ast$ mide el beneficio.}

% ==========================================================================
\section{Cómo se calcula el óptimo}

Conviene saberlo, porque el applet no hace lo que hace el original de GeoGebra
y la diferencia se nota en los casos interesantes.

\textbf{No se resuelve $CMg = p$.} Esa ecuación tiene dos raíces en una curva de
costos cúbica, y una de ellas es un mínimo del beneficio. Con cualquier función
de costos no convexa el óptimo global puede estar además en una esquina, donde
ninguna condición de primer orden lo encuentra.

Lo que se hace es maximizar el beneficio directamente:

\begin{enumerate}
\item Se acota el rango por arriba, buscando el punto a partir del cual el
  costo marginal ha superado el precio y sigue subiendo.
\item Se barre ese rango con 700 puntos y se queda con el mejor.
\item Se refina ese entorno por sección áurea.
\item Se compara el resultado con cerrar, que da $\pi(0) = -F$. Cerrar es un
  candidato de pleno derecho, no un caso límite.
\end{enumerate}

La curva de oferta se traza igual: una maximización completa por cada uno de
140 precios. Es la única manera de que salga el salto. Invirtiendo $CMg = p$
nunca aparecería, porque por debajo del precio de cierre la respuesta es cero y
en el precio de cierre da un brinco.

\begin{amnote}
El mínimo del costo variable medio se busca sobre una escalera geométrica, no
lineal: lo que decide si la empresa llega a cerrar alguna vez es la forma de la
curva cerca de cero, y una malla uniforme sobre un rango de cinco órdenes de
magnitud no resuelve nada ahí. Si el mínimo cae en un extremo de la escalera es
que no hay mínimo —la curva sigue bajando— y el applet informa de que no hay
precio de cierre en vez de dar el valor más pequeño que muestreó.
\end{amnote}

% ==========================================================================
\section{El veredicto}

La banda de arriba nombra el caso y lo explica. Son seis.

\begin{ctrltable}{Veredicto & Cuándo}
Produce con beneficio & $p > p_{ne}$. El precio supera el costo medio mínimo. \\
Produce con pérdidas & $\underline{p} < p < p_{ne}$. Y hace bien: cerrar
  costaría el costo fijo entero. \\
Beneficio nulo & $p = p_{ne}$, el mínimo del costo medio. \\
Cierra & $p < \underline{p}$. El óptimo es $q^\ast = 0$ y la pérdida es
  exactamente $F$. \\
Contra la capacidad & El óptimo está prácticamente en $\bar Q$. Solo aparece en
  la familia logística. \\
No hay óptimo & El costo marginal no supera al precio en ningún rango
  dibujable. \\
\end{ctrltable}

\fig[1.0]{shutdown}{Cierre, con la cúbica a $p = 1{,}15$. La curva de oferta
está pegada al eje de precios, $q^\ast$ se lee en rojo, y el veredicto explica
por qué producir sería peor que no producir.}

% ==========================================================================
\section{Actividades para clase}

Cada una empieza desde los ajustes de partida, salvo que se diga otra cosa. Si
se ha perdido, recargue la página.

\begin{activity}{Las dos curvas son una}
\amset Función de costos \ui{Cuadráticos / potencia}, valores de partida,
$p = 2$. Encienda \ui{Costo marginal} si estuviera apagado.
\amexp La curva cian del panel izquierdo y la roja del derecho ocupan el mismo
lugar en la pantalla. Mueva $p$: el punto morado se desplaza por las dos a la
vez, siempre a la misma altura. Suba $a$ o $k$ y las dos se empinan juntas.
\par\smallskip
Pregunte a la clase por qué el panel izquierdo no tiene curva de costo medio
dibujada con trazo grueso: porque el costo medio no interviene en la decisión
de cuánto producir, solo en la de si el negocio merece la pena.
\end{activity}

\begin{activity}{Fabricar un precio de cierre}
\amset Elija \ui{Personalizada}. La expresión por defecto es la cúbica.
\amexp Ahora $CVMe$ tiene forma de U con mínimo $1{,}375$, y la oferta ya no
sale del origen: hay un tramo pegado al eje de precios y un salto.
\par\smallskip
Baje el precio hasta $1{,}30$ y luego súbalo a $1{,}40$. La cantidad pasa de $0$
a $2{,}55$ de golpe. Pregunte: ¿qué cantidad produce la empresa cuando
$p = 1{,}375$ exactamente? La respuesta honesta es que le da igual, y por eso
la oferta ahí no es una función sino una correspondencia.
\end{activity}

\begin{activity}{Producir perdiendo dinero}
\amset Cúbica. En esta cúbica $F = 0$, así que el precio de cierre y el de
nivelación coinciden en $1{,}375$ y no hay franja de pérdidas. Añádale un costo
fijo: escriba \code{1+0.1q\^{}3-0.5q\^{}2+2q} en la casilla. Ponga $p = 1{,}6$.
\amexp El precio de cierre sigue siendo $1{,}375$ —el costo fijo no toca al
costo variable medio— pero el de nivelación sube a $1{,}733$. Entre los dos hay
ahora una franja de precios, y $1{,}6$ está dentro de ella.
\par\smallskip
La lectura \ui{Beneficio} da $-0{,}394$ en rojo, y aun así el óptimo es producir
$2{,}869$. Pregunte por qué no cierra. Cerrar costaría el costo fijo entero,
$-1$; produciendo pierde solo $0{,}394$. Mueva el precio por toda la franja
$1{,}375 < p < 1{,}733$: el beneficio es negativo en toda ella y producir sigue
siendo lo mejor.
\end{activity}

\begin{activity}{El excedente medido dos veces}
\amset Cúbica, $p = 1{,}6$. Encienda \ui{Área de excedente}.
\amexp Aparecen dos regiones sombreadas de formas claramente distintas, una en
cada panel. La lectura \ui{Excedente} da $0{,}606$, y es la de las dos.
\par\smallskip
Compruébelo con los números del panel derecho: $q^\ast = 2{,}869$,
$\text{ingreso} = p\,q^\ast = 4{,}59$, $c(q^\ast) = 3{,}98$. Como aquí $F = 0$,
el costo variable es el costo entero, y
$EP = 4{,}59 - 3{,}98 = 0{,}61$. Coincide con las dos áreas.
\par\smallskip
Pregunte por qué las dos regiones no tienen la misma forma si miden lo mismo.
Porque una integra $p - CMg$ respecto de la cantidad y la otra integra $q^\ast$
respecto del precio; con un salto en el precio de cierre no pueden ser la misma
región, pero el teorema fundamental del cálculo dice que tienen la misma área.
\end{activity}

\begin{activity}{La condición de primer orden, dibujada}
\amset Cúbica, $p = 1{,}6$, panel de la empresa en \ui{Totales}.
\amexp La tangente punteada al costo total en $q^\ast$ es paralela a la recta
de ingresos. Mueva el precio: la recta de ingresos gira, y el punto de
tangencia se desliza por la curva de costos hasta donde la pendiente vuelve a
coincidir.
\par\smallskip
El segmento vertical en $q^\ast$ es el beneficio. Encienda \ui{Punto de
equilibrio} y aparecen las dos cantidades en las que ese segmento tiene
longitud cero.
\end{activity}

\begin{activity}{Dos raíces, un óptimo}
\amset Cúbica, $p = 1{,}6$, vista \ui{Marginales}.
\amexp La horizontal del precio corta al costo marginal \emph{dos} veces. El
punto morado está en la de la derecha, donde $CMg$ está subiendo.
\par\smallskip
Pregunte qué pasa en la otra. Es la solución de $CMg = p$ que un estudiante
obtendría resolviendo la ecuación sin mirar la segunda derivada, y es un mínimo
del beneficio: producir esa cantidad es lo peor que puede hacer la empresa
entre las que le costean el variable. Pásese a \ui{Totales} para verlo como un
valle en la curva de beneficio.
\end{activity}

\begin{activity}{Capacidad}
\amset Familia \ui{Producción log-CD logística}, valores de partida.
\amexp La vertical gris marca $\bar Q = b\bar K^{\gamma}$. Suba el precio hasta
el tope del deslizador: la cantidad crece cada vez menos, y la lectura
\ui{Elasticidad} se acerca a cero.
\par\smallskip
Suba ahora $\bar K$ y la capacidad se desplaza a la derecha. Ese es exactamente
el ejercicio de largo plazo que el corto plazo no permite: la única manera de
producir más allá de $\bar Q$ es cambiar de planta.
\end{activity}

\begin{activity}{Cuando no hay óptimo}
\amset Familia \ui{Producción Cobb--Douglas}. Baje $\beta$ a $0{,}05$ y suba
$s$ a $2$.
\amexp La fórmula de la columna muestra $e = 0{,}526$, y el veredicto pasa a
«No hay óptimo». Todas las lecturas del panel de oferta se vuelven rayas.
\par\smallskip
El costo marginal es decreciente: cada unidad extra cuesta menos que la
anterior y se vende al mismo precio. Pregunte qué le pasaría a una empresa así
en un mercado competitivo. La respuesta —que crecería hasta dejar de ser
competitiva— es el argumento por el que los rendimientos decrecientes no son un
supuesto técnico sino la condición para que el modelo tenga sentido.
\end{activity}

% ==========================================================================
\section{Sintaxis de expresiones}

En la casilla de costos la variable es \code{q}. También se acepta \code{x},
porque es como se llamaba en GeoGebra.

\begin{ctrltable}{Elemento & Qué se admite}
Operadores & \code{+ - * / \^{} ( )} \\
Funciones & \code{sqrt ln log log10 exp abs sin cos tan min max pow} \\
Constantes & \code{e}, \code{pi} \\
Multiplicación implícita & \code{2q}, \code{3(q+1)}, \code{2sqrt(q)} \\
Asociatividad & \code{\^{}} liga por la derecha y más fuerte que el menos
  unario: \code{-q\^{}2} es $-(q^2)$ \\
\end{ctrltable}

El costo marginal se deriva simbólicamente. Cuando la expresión tiene un
quiebre —\code{min}, \code{max}, \code{abs}— no hay derivada simbólica que
valga y se pasa a una diferencia centrada, que es lo que se ve dibujado de
todos modos.

\begin{amwatch}
Escribir una función de costos decreciente, o con costo marginal constante,
lleva al caso «No hay óptimo». No es un fallo del applet: es que ese problema
no tiene solución.
\end{amwatch}

% ==========================================================================
\section{Diferencias con el original de GeoGebra}

\begin{description}
\item[Dos paneles en lugar de uno.] El original dibujaba costo total, ingreso,
  beneficio, costo marginal y costo medio sobre los mismos ejes. El dinero y el
  dinero por unidad no comparten escala, así que las curvas marginales quedaban
  ilegibles al lado de las totales. Aquí están separadas, y el panel marginal
  está enganchado a la escala del de oferta.
\item[La curva de oferta, que no estaba.] Junto con el precio de cierre, el de
  nivelación, el salto entre ellos y la elasticidad.
\item[El óptimo es una maximización, no una intersección.] El original
  intersecaba \code{Eq: CM(x) = p} con el eje y usaba el resultado como extremo
  derecho de un \code{Max[Beneficios, 0, x(A)]}. Con costo marginal no monótono
  ese intervalo es el equivocado, y cuando $CMg = p$ no tiene raíz la cadena
  entera devolvía \code{NaN} y el applet se quedaba en blanco.
\item[Cerrar es un candidato.] El original siempre comparaba contra un máximo
  interior. Aquí $\pi(0) = -F$ compite con él, que es lo que hace que el precio
  de cierre exista.
\item[Un ínfimo no es un mínimo.] Un costo variable medio que solo decrece
  hacia cero no tiene mínimo, luego no hay precio de cierre. Informar del valor
  más pequeño muestreado daba un precio de cierre de $7{,}7\times10^{-4}$.
\item[Deslizadores acotados lejos de la división por cero.] $\beta = 1$ divide
  por $1-\beta$; $\bar K = 0$ eleva cero a una potencia negativa; $w = 0$ y
  $a = 0$ dejan el costo constante; $k = 1$ deja el costo marginal constante.
  Todos eran alcanzables en el original.
\item[$m$ ampliado] de $[0, 1]$ a $[-1, 2]$, donde de verdad mueve el nivel de
  costos lo bastante para verse.
\item[Bilingüe, y con veredicto.] Cada estado del problema se nombra y se
  explica, porque las formas por sí solas no dicen en cuál de los seis casos
  está mirando el estudiante.
\item[Objetos muertos no trasladados.] \code{AreaProfi}, \code{CM2} (una copia
  de \code{CM}), \code{eq1}, \code{A}, \code{prueba} y \code{prueba2} eran
  fontanería de la construcción de GeoGebra, no objetos con significado.
\end{description}

% ==========================================================================
\section{Problemas frecuentes}

\begin{ctrltable}{Síntoma & Qué pasa}
Los paneles salen vacíos & La expresión de costos no se entiende. El mensaje en
  rojo bajo la casilla dice si es un problema de sintaxis o de variable. \\
Todo dice «$\infty$» o rayas & Está en el caso «No hay óptimo». Suba $k$, o
  $s$, o $\beta$, o cambie de familia. \\
Las escalas dejan de coincidir & El panel de la empresa está en \ui{Totales}.
  La pastilla de su cabecera lo dice. Vuelva a \ui{Marginales}. \\
El eje se mueve al arrastrar el precio & Es deliberado: los ejes se redondean a
  números redondos y saltan por escalones, para que dos ajustes puedan
  compararse a ojo en vez de renumerarse continuamente. \\
El precio no responde al ratón & Está arrastrando dentro del panel de totales,
  cuyo eje vertical es dinero. Use el panel de oferta o el deslizador. \\
Se perdió el idioma o el tema & Se guardan en el navegador. En una ventana
  privada, o tras borrar los datos del sitio, vuelven a sus valores de
  partida. \\
\end{ctrltable}

%% Secciones comunes a todos los manuales, edición española.
%% Se incluyen con \input{common-es}. Lo único que cambia entre applets es
%% \appletfolder, que es también el nombre de la carpeta y de la ruta en el sitio.

\section{Publicarlo para una clase}

Sirve cualquier alojamiento estático: GitHub Pages, Netlify, un directorio web
de la universidad, o incluso una carpeta compartida servida por HTTP. Lo único
que hay que subir es el archivo \code{index.html} dentro de una carpeta con el
nombre que quiera que aparezca en la dirección.

El repositorio trae un flujo de trabajo de GitHub Actions
(\code{.github/workflows/pages.yml}) que publica todos los applets a la vez en
cada empuje a la rama principal. Para activarlo, en el repositorio:
\emph{Settings\,$\rightarrow$\,Pages\,$\rightarrow$\,Source: GitHub Actions}.
Es lo único que no puede hacer el propio flujo, porque crear un sitio de Pages
requiere permisos de administración que un \code{GITHUB\_TOKEN} nunca recibe.

\begin{amnote}
Para repartir el applet a estudiantes sin conexión fiable, mándeles el archivo
\code{index.html} directamente. Pesa menos que una fotografía y se abre con
doble clic.
\end{amnote}

\section{Licencia y procedencia}

El código de la aplicación es MIT. Puede usarlo, modificarlo y repartirlo, en
clase o fuera de ella, citando la procedencia.

Este applet es una reescritura autónoma de un applet original de GeoGebra. La
reescritura no es una traducción literal: donde el original tenía un error de
construcción, una división por cero alcanzable con los deslizadores, o un objeto
muerto heredado del applet del que se copió, esta versión lo corrige y lo
documenta. La sección \emph{Diferencias con el original de GeoGebra} enumera
esos cambios uno por uno.

Todo está escrito a mano y sin dependencias: el analizador de expresiones, la
derivación simbólica, el trazado de curvas de nivel y el dibujo. En particular
no se usa \code{eval}, de modo que lo que escribe un estudiante en una casilla
de texto no puede convertirse en código, y una política de seguridad de
contenidos estricta no rompe la página.


\end{document}
